\documentclass{beamer}
\usepackage{beamerthemesplit}

\usetheme{SPbGU}

\usepackage[T2A]{fontenc} 
\usepackage[utf8]{inputenc}
\usepackage[english,russian]{babel}

\usepackage{wrapfig}
\usepackage{graphicx}
\usepackage{caption}
\usepackage{subcaption}
\usepackage{pdfpages}
\usepackage{multirow}
\usepackage{fancyvrb}

\usepackage{soulutf8}
\usepackage{amsmath}
\usepackage{cmap} 
\usepackage{indentfirst}
\usepackage{amsmath}
\usepackage{tikz}
\usepackage{textcomp}

\usepackage[noend]{algpseudocode}
\usepackage{algorithm}
\usepackage{algorithmicx}
\usepackage{listings}

\newtheorem{rutheorem}{Теорема}
\newtheorem{ruproof}{Доказательство}
\newtheorem{rudefinition}{Определение}
\newtheorem{rulemma}{Лемма}

\usetikzlibrary{shapes,arrows}

\beamertemplatenavigationsymbolsempty

\title[]{Динамически формируемый код: синтаксический анализ контекстно-свободной аппроксимации}
\subtitle[]{}
% То, что в квадратных скобках, отображается в левом нижнем углу. 
\institute[СПбГУ]{
	Санкт-Петербургский государственный университет\\
	Лаборатория языковых инструментов JetBrains }

% То, что в квадратных скобках, отображается в левом нижнем углу.
\author[Дмитрий Ковалев]{Дмитрий Ковалев}

\date{4 апреля 2017г.}

\definecolor{orange}{RGB}{179,36,31}

\begin{document}
{

 %Лого университета или организации, отображается в шапке титульного листа
\begin{frame}

\begin{figure}
	\centering
	\begin{subfigure}[h]{0.3\textwidth}
		\centering
		\includegraphics[width=1.5cm]{pictures/SPbGU_Logo.png}
	\end{subfigure}
	\begin{subfigure}[h]{0.3\textwidth}
		\centering
		\includegraphics[width=3.5cm]{pictures/jetbrains.pdf}
	\end{subfigure}
\end{figure}
\titlepage
\end{frame}
}
%\maketitle

\begin{frame}[fragile]
	\transwipe[direction=90]
	\frametitle{Динамически формируемый код}
	\begin{itemize}
		\item SQL-запросы в C$\#$
			\begin{Verbatim}[commandchars=\\\{\}]
  \textcolor{blue}{private void} Example (\textcolor{blue}{bool} cond) \{
      \textcolor{blue}{string} columnName = cond ? \textcolor{red}{"name"} : \textcolor{red}{"address"};
      \textcolor{blue}{string} queryString = 
          \textcolor{red}{"SELECT id, "} + columnName + \textcolor{red}{" FROM users"};
      Program.ExecuteImmediate(queryString);
  \}
			\end{Verbatim}
		\pause
		\item Генерация HTML-страниц в PHP
			\begin{Verbatim}[commandchars=\\\{\}]	
  <?php 
      \$name = \textcolor{red}{'your name'};
      echo \textcolor{red}{'<table>} 
           \textcolor{red}{<tr><th>Name</th></tr>}  
           \textcolor{red}{<tr><td>'}.\$name.\textcolor{red}{'</td></tr>} 
           \textcolor{red}{</table>'};
  ?>
			\end{Verbatim}
	\end{itemize}
\end{frame}

\begin{frame}[fragile]
	\transwipe[direction=90]
	\frametitle{Статический анализ динамически формируемого кода} 
	\begin{itemize}
		\item Необходимо анализировать каждое возможное значение выражения
		\item В общем случае задача неразрешима 
		\item Для анализа используется аппроксимация множества значений
	\end{itemize}
	
	\pause
	\begin{itemize}
		\item Общая схема:
		\begin{itemize}
			\item построение аппроксимации
			\item лексический анализ
			\item \textbf{синтаксический анализ}
			\item семантический анализ
		\end{itemize}
	\end{itemize}
\end{frame}

\begin{frame}[fragile]
	\transwipe[direction=90]
	\frametitle{Синтаксический анализ регулярной аппроксимации}
	\begin{figure}
		\centering
		\includegraphics[width=12cm]{pictures/scheme_full.pdf}
	\end{figure}
	\pause
	\vspace{-15pt}
	\only<2>{\begin{itemize}
			\item Проверка включения регулярного языка в LL(k)
			\item Лес разбора --- множество деревьев разбора корректных строк
		\end{itemize}}
	\only<3>{\begin{itemize}
			\item Существующие решения основаны на ...
			\begin{itemize}
				\item GLR-алгоритме (Alvor, 2010)
				\item RNGLR (Е. Вербицкая, 2015)
				\item \textbf{GLL} (А. Рагозина, 2016)
			\end{itemize}
		\end{itemize}}
\end{frame}

\begin{frame}[fragile]
	\transwipe[direction=90]
	\frametitle{Generalized LL}
	\begin{itemize}
		\item Работает с произвольными КС-грамматиками
		\item Результат --- Shared Packed Parse Forest (SPPF)
		\item Структурированный в виде графа стек (GSS)
		\begin{itemize}
			\item Экономия памяти
			\item Гарантирует завершаемость алгоритма
		\end{itemize}
	\end{itemize}
	\begin{figure}
		\centering
		\includegraphics[width=9cm]{pictures/gss.pdf}
	\end{figure}
\end{frame}

\begin{frame}[fragile]
	\transwipe[direction=90]
	\frametitle{GLL и нелинейный вход}
	\begin{itemize}
		\item Механизм дескрипторов --- $(L, u, i, w)$, где
		\begin{itemize}
			\item $L$ --- позиция в грамматике вида $A \rightarrow \alpha \cdot X \beta$
			\item $u$ --- вершина GSS
			\item $i$ --- \only<1>{позиция во входной строке}
										\only<2>{\st{позиция во входной строке} \ номер вершины автомата}
			\item $w$ --- узел SPPF
		\end{itemize}
		\pause
		\item Вместо следующего символа просматриваются метки всех исходящих ребер
	\end{itemize}
\end{frame}

\begin{frame}[fragile]
	\transwipe[direction=90]
	\frametitle{Неточность регулярной аппроксимации}
	\begin{table}[h]
		\centering
		\begin{tabular}{p{5cm} p{5cm}}
			\begin{center}
				\begin{Verbatim}[commandchars=\\\{\}]
expr = \textcolor{red}{"a"}
while ...
    expr = \textcolor{red}{"("} + expr + \textcolor{red}{")"}
print expr
				\end{Verbatim}
			\end{center}
			&
			\vspace{-15pt}
			\begin{center}
				\includegraphics[width=5cm]{pictures/reg_app.pdf}	
			\end{center} \\
			Порождаемый язык: $(^n \; a \; )^n$
			&
			$\quad$ Порождаемый язык: $(^* \; a \; )^*$
		\end{tabular}
	\end{table}
	\begin{itemize}
		\item Сообщения об ошибках, отсутствующих в строках из исходного множества значений
	\end{itemize}
\end{frame}

\begin{frame}[fragile]
	\transwipe[direction=90]
	\frametitle{Синтаксический анализ КС-аппроксимации}
	\begin{figure}
		\centering
		\includegraphics[width=12cm]{pictures/scheme_cs.pdf}
		\begin{itemize}
			\item Получение КС-аппроксимации множества значений выражения $\,$ (Y. Minamide, 2005)
		\end{itemize}
	\end{figure}
\end{frame}

\begin{frame}[fragile]
	\transwipe[direction=90]
	\frametitle{Разрешимость}
	\begin{itemize}
		\item Проверка включения КС-языков --- неразрешимая задача
		\item Класс КС-языков не замкнут относительно пересечения
	\end{itemize}
	
	\pause
	\begin{itemize}
		\item Для некоторых подклассов задача разрешима
	\end{itemize}
	\begin{itemize}
		\item Приближенные решения:
		\begin{itemize}
			\item Doh, K. G. Abstract Parsing: Static Analysis of Dynamically \\ Generated String Output Using LR-Parsing Technology
		\end{itemize}
	\end{itemize}
\end{frame}

\begin{frame}[fragile]
	\transwipe[direction=90]
	\frametitle{Графовое представление грамматики}
	\vspace{-20pt}
	\begin{table}
		\centering
		\begin{tabular}{p{5cm} p{7cm}}
			\begin{center}
				\hspace{30pt}
				\begin{BVerbatim}[commandchars=\\\{\}]
s\textquotesingle ::= s
s  ::= LBR s RBR | A
				\end{BVerbatim}
			\end{center}
			&
			\begin{center}
				\multirow{-3}*{\includegraphics[width=4cm]{pictures/RA.pdf}}
			\end{center}
		\end{tabular}
	\end{table}
	\vspace{70pt}
	\begin{itemize}
		\item Ребро с нетерминалом --- переход в стартовую вершину данного нетерминала
	\end{itemize}
\end{frame}

\begin{frame}[fragile]
	\transwipe[direction=90]
	\frametitle{Модификация алгоритма}
	\begin{itemize}
		\item Необходимо поддерживать два стека (GSS) 
		\pause
			\item Расширенные дескрипторы --- $(L, u_1, i, u_2, w)$, где
			\begin{itemize}
				\item $L$ --- позиция в грамматике вида $A \rightarrow \alpha \cdot X \beta$
				\item $u_1$ --- вершина первого GSS
				\item $i$ --- номер вершины автомата (позиция во второй грамматике)
				\item $u_2$ --- вершина второго GSS
				\item $w$ --- узел SPPF
			\end{itemize}
		\end{itemize}
\end{frame}

\begin{frame}[fragile]
	\transwipe[direction=90]
	\frametitle{Результат работы алгоритма}
	\begin{table}[h]
		\centering
		\begin{tabular}{p{5cm} p{5cm}}
			\begin{center}
				\begin{Verbatim}[commandchars=\\\{\}]
expr = \textcolor{red}{"a"}
while ...
    expr = \textcolor{red}{"("} + expr + \textcolor{red}{")"}
print expr
				\end{Verbatim}
			\end{center}
			&
			\vspace{-15pt}
			\begin{center}
				\includegraphics[width=5cm]{pictures/sppf.pdf}	
			\end{center}
	\end{tabular}
\end{table} 
\vspace{-20pt}
\begin{table}
	\centering
	\begin{tabular}{p{2.5cm} | c}
		Аппроксимация & Кол-во ошибок \\ \hline
		Регулярная & 7 \\ \hline
		\quad \quad КС & 0
	\end{tabular}
\end{table}
\end{frame}

\begin{frame}[fragile]
	\transwipe[direction=90]
	\frametitle{Заключение}
	\begin{itemize}
		\item Представлен алгоритм анализа КС-аппроксимации динамически формируемого кода на основе GLL
		\item Алгоритм реализован в рамках проекта YaccConstructor
	\end{itemize}
	
\end{frame}


\end{document}
